\documentclass[11pt]{article}
\usepackage[margin=1.15in]{geometry}
\usepackage{booktabs}
\usepackage{amsmath,amssymb}
\usepackage{microtype}
\usepackage[numbers,sort&compress]{natbib}
\usepackage[hidelinks]{hyperref}
\usepackage{xcolor}
\usepackage{enumitem}
\usepackage{titlesec}

\titleformat{\section}{\normalfont\large\bfseries}{\thesection}{0.6em}{}
\titleformat{\subsection}{\normalfont\normalsize\bfseries}{\thesubsection}{0.6em}{}
\setlist{itemsep=2pt,topsep=4pt,parsep=0pt}

\newcommand{\cost}{\ensuremath{\mathrm{FN}\cdot c + \mathrm{FP}}}

\title{\bfseries What Is a False Negative Worth?\\[2pt]
\large Grounding Asymmetric Cost in Labor Hours, and the Threshold Bug It Uncovers}

\author{Kartik Thakore\\\small kartik@multiversal.ventures}
\date{\today}

\begin{document}
\maketitle

\begin{abstract}
\noindent
Cost-sensitive learning asks practitioners for one number: how many false alarms a missed
positive is worth. In applied work that number is almost always asserted. We show what
asserting it costs, on a corpus of 156{,}732 SEC filings labelled for later retraction, and
we give a procedure for deriving it instead from published labor rates.

Three findings. First, the ratio decides the result: three independent feature spaces agree
that our models beat the trivial baseline below $c \approx 5.6$ and lose above it, and the
constant our own work had inherited, $c = 10$, turned out to be a model-selection key
borrowed from an unrelated task in another repository. Second, deriving the ratio from the
measured audit-fee premium of a restatement puts it between 24 and 300, far above the
crossover, but also reveals that the ratio is a property of \emph{who is asking} rather than
of the task: a preparer, an investor and a regulator have three different numerators, and one
of them is not a ratio at all. Third, while measuring this we found that the ratio was not
even the largest error in our own pipeline. Selecting the decision threshold on in-sample
scores inflates test cost by up to $5.5\times$. The inflation is governed by how completely
the model memorises its training set, not by $p/n$ or by model quality, so it is largest for
the models most likely to be deployed. Calibration provably cannot repair it. Out-of-fold
selection recovers a median $95\%$ of the loss at the cost of $K$ extra fits and no held-out
data.

We close with a checklist. The practical recommendation is to report a cost sweep rather than
a point, to derive any point estimate from hours somebody can be asked about, and where the
seat is ambiguous to evaluate at a stated review capacity instead, since $\mathrm{precision}@k$
requires no cost ratio at all.
\end{abstract}

\section{Introduction}

A practitioner building a detector for a business process is handed an asymmetric problem. A
missed fraud, a missed defect, a missed restatement is worse than a false alarm. The standard
response is a cost-sensitive loss,
\[
\text{cost} \;=\; \mathrm{FN}\cdot c \;+\; \mathrm{FP},
\]
and the standard practice is to pick $c$ by assertion: a round number, a number from a paper
in the same domain, a number a stakeholder offered in a meeting.

This paper is about what that costs. It is written for practitioners, and every claim in it is
measured on real data with code and seeds fixed.

\paragraph{Contributions.}
\begin{enumerate}[label=(\arabic*)]
\item \textbf{We show the ratio decides the answer, and quantify where the decision flips.} On
our corpus the crossover is $c \in [5.5, 5.75]$ across three unrelated feature spaces. Below
it every model beats the trivial baseline; above it every model loses. Nothing about the
models changes.
\item \textbf{We give a procedure for grounding $c$ in labor hours} and run it. This is, to our
knowledge, the first attempt in this literature to derive an asymmetric cost ratio from
published compensation and fee data rather than assert it.
\item \textbf{We identify a failure of that procedure that is not a data gap.} The derived
ratio depends on the seat the model is deployed in, and a task specification that does not
name a seat does not determine a ratio. We argue this is common and under-recognised.
\item \textbf{We report a threshold-selection artefact} worth up to $5.5\times$ in test cost,
establish its mechanism, show that the obvious remedy provably cannot work, and benchmark eight
remedies that can.
\end{enumerate}

\section{The ratio is asserted, not measured}
\label{sec:asserted}

Cost-sensitive learning has a mature theory \citep{elkan2001,provost2001} and mature machinery:
MetaCost, class weighting, instance-dependent variants \citep{bahnsen2014}. The theory takes
the cost matrix as given. The applied literature then supplies it by assertion. Surveying
recent financial-misstatement work, we found ratios of $6.46{:}1$ and similar figures stated
without derivation.

Our own case is worse and is instructive. Our repository's constant traced to a single table
row in a different repository:

\begin{quote}\small
\texttt{| selection | val flip-pair accuracy | val asymmetric key (missed$\times$10 + false) |}
\end{quote}

It was the metric used to select a training epoch, for a different model, on 178 synthesised
rows, answering a different question. It was never a claim about our domain. Every result we
had published was reported against it as though it described the world.

An unmeasured constant is at least a guess about the right domain. This one had crossed a
repository, a corpus and a decision problem. We suspect this is not rare.

\section{The ratio decides the result}
\label{sec:sweep}

\paragraph{Setup.} 156{,}732 SEC filings from 5{,}085 companies, 2004--2024. Labels come from
8-K Item 4.02, the item a registrant files to state that previously issued financials should no
longer be relied upon; the 8-K names the filings it retracts, so labels are mechanical rather
than annotated. Two views per filing: narrative section text, and XBRL facts. Splits are by
company throughout, never by row, because filings from one company are not independent.

We score each arm against \emph{the best floor}, which is the cheaper of always-fire
($\text{cost} = \mathrm{N}$) and never-fire ($\text{cost} = \mathrm{P}\cdot c$) at that ratio.
Which floor wins is itself a function of $c$, so comparing against a fixed one confounds two
effects.

\begin{table}[t]
\centering\small
\begin{tabular}{lccc}
\toprule
$c$ & XBRL arm & narrative arm & embedding arm \\
\midrule
3  & 0.652 & 0.576 & 0.661 \\
5  & 0.903 & 0.882 & 0.916 \\
10 & 1.572 & 1.407 & 1.273 \\
\bottomrule
\end{tabular}
\caption{Arm cost over the best floor. Below 1.0 the model wins. Same models, same rows, same
split; only $c$ moves. On a quarter-point grid the crossover is 5.5, 5.75 and 5.5.}
\label{tab:sweep}
\end{table}

Table~\ref{tab:sweep} is the whole argument of Section~\ref{sec:asserted} in one place. Three
feature spaces with nothing in common agree on the crossover to within a quarter of a point,
and the constant we had inherited sits on the far side of it.

\subsection{What the linear form hides}

The form itself has a defect independent of the constant. Under \cost{}, always-fire reviews
$100\%$ of cases and pays only the per-review price. It is therefore purchasable at any scale,
and it wins whenever the base rate clears $1/(c+1)$.

\begin{quote}
\emph{No ranking model can beat a strategy that takes every positive by taking every row.}
\end{quote}

A corpus scored this way is answering ``is the base rate high'' while appearing to answer ``does
the model discriminate.'' Writing the cost per row with base rate $\pi$,
\[
\frac{\text{cost}}{n} \;=\; c\,\pi\,(1-\mathrm{TPR}) \;+\; (1-\pi)\,\mathrm{FPR},
\]
always-fire is the ROC point $(1,1)$, and a model beats it exactly when
\[
c \;<\; \underbrace{\frac{1-\pi}{\pi}}_{\text{odds of a negative}}\;\cdot\;
\underbrace{\max_{\mathrm{ROC}} \frac{1-\mathrm{FPR}}{1-\mathrm{TPR}}}_{\text{base-rate free}} .
\]
This is a restatement of iso-performance analysis \citep{provost2001} and is not new geometry.
What is worth saying to practitioners is the consequence: the contest is decided by the steepest
line from the \emph{top-right} corner down to the ROC curve, and a model can have excellent AUC
while being flat exactly there. Reading a ROC curve from $(0,0)$ tells you nothing about it.

\section{Grounding the ratio in labor hours}
\label{sec:hours}

If the ratio must be supplied, supply it from something a person can be asked about. Both halves
of the following are answerable by an operations lead:
\[
c \;=\; \frac{\text{hours a miss eventually costs}}{\text{hours one review costs}} .
\]

\paragraph{The recipe.}
\begin{enumerate}[label=\textbf{\arabic*.}]
\item \textbf{Name the seat.} Who acts on the prediction, and what do they do differently when
it fires? This step is not optional; Section~\ref{sec:seat} is what happens when it is skipped.
\item \textbf{Price the review.} Hours per case at that seat, times a loaded rate. Triage and
full review differ by an order of magnitude and the answer moves with them.
\item \textbf{Price the miss, in the same units.} Published fee premia, remediation studies,
incident post-mortems, and internal time-tracking are all admissible. State which.
\item \textbf{Apply the detection discount.} A flagged case only avoids the downstream cost if
the review actually finds the problem. Multiply by that rate, call it $d$, and be explicit that
you are guessing at it if you are.
\item \textbf{Report the grid, not the point.} Steps 2 to 4 each carry uncertainty. A table
across their plausible ranges is a result; a single number is a claim you cannot support.
\end{enumerate}

\paragraph{Worked example.} An interim restatement carries a mean audit fee premium of
\$191{,}000, the gap between a restating client's \$1.583M mean fee and the \$1.392M it would
otherwise have paid \citep{interim2025}. At the published US blended range of \$150 to \$450 per
hour that is 424 to 1{,}273 professional hours. It is the auditor alone; internal controller,
legal and consultant hours do not appear in any standardised disclosure, so this is a lower
bound rather than an estimate. Taking $H_\text{miss} = 760$ hours (the midpoint rate):

\begin{table}[t]
\centering\small
\begin{tabular}{lrrrrr}
\toprule
detection rate $d$ & 0.25\,h review & 0.5\,h & 1\,h & 4\,h & 8\,h \\
\midrule
0.10 & 304  & 152  & 76  & 19 & \textbf{10} \\
0.25 & 760  & 380  & 190 & 48 & 24 \\
0.50 & 1520 & 760  & 380 & 95 & 48 \\
1.00 & 3040 & 1520 & 760 & 190 & 95 \\
\bottomrule
\end{tabular}
\caption{Implied $c$, from $c = d \cdot H_\text{miss} / H_\text{review}$ with
$H_\text{miss}=760$ professional hours. One cell in twenty reaches the constant our work had
been using, and it is the corner: a review that finds the problem one time in ten, costing a
full working day.}
\label{tab:hours}
\end{table}

Every cell but one is between two and three hundred times the asserted constant, and the
crossover from Table~\ref{tab:sweep} is 5.5. The practical reading is uncomfortable and worth
stating plainly: \textbf{at the ratio the domain actually implies, flagging everything is
correct on cost, and no model in this study was ever going to win that argument.}

\subsection{The seat problem}
\label{sec:seat}

The 760 hours belongs to a seat our corpus does not occupy, and noticing this is the most
transferable result in the paper.

Our task is post-filing: given a document already filed, predict whether it will later be
retracted. Flagging it \emph{does not prevent the restatement}. The error is in the numbers and
the remediation happens either way. What early detection buys depends on who is asking:

\begin{itemize}
\item \textbf{A preparer or auditor, before filing.} Catching it here does avoid the cost, and
760 hours is the right numerator. But this seat cannot be our corpus, because the label only
exists after the fact.
\item \textbf{An investor.} A miss is holding a position through an announcement effect of
roughly $-9.2\%$ \citep{gao2002}. The cost scales with position size, so a unitless ratio
against ``one analyst hour'' is not the right shape for the question.
\item \textbf{A regulator or reviewer triaging filings.} The remediation happens regardless.
What is bought is \emph{earlier}, valued in months rather than in hours avoided.
\end{itemize}

\begin{quote}
\emph{A task specification that does not name a seat does not determine a cost ratio.}
\end{quote}

This is not a gap that better data fills. It is a statement about what the task is. We suspect
many benchmark tasks with asserted cost matrices are in this position, and that the asserted
constant is doing the work of a missing problem statement.

\paragraph{The way out.} Where the seat is genuinely ambiguous, evaluate at a stated
\emph{capacity} instead. A reviewer has $k$ slots, so the decision is a queue and not a label,
and always-fire leaves the choice set because it cannot say which $k$ to open.
$\mathrm{precision}@k$ requires no cost ratio at all \citep{consequentialist2025}, only a number
every operations team already knows.

The change is not cosmetic. Under a top-$1\%$ budget on the 2021 window, our narrative arm's
top 11 ranked filings are all later-retracted, and so are its top 22: it reaches the $1/\pi$
ceiling exactly. Under $\mathrm{FN}\cdot 10 + \mathrm{FP}$ the same arm on the same scores reads
as $41\%$ worse than useless. Both statements are true, and only one describes something a team
could staff.

\section{The threshold was a larger error than the ratio}
\label{sec:threshold}

While measuring the above we found a defect in our own pipeline that dominated it, and that we
believe is widespread because it is invisible.

A score is not a decision. Somewhere a threshold is chosen. Our code, and much applied code,
chose it by sweeping the model's scores \emph{on the rows the model was fitted on}.

\subsection{Measurement}

Table~\ref{tab:strategies} benchmarks eight threshold-placement strategies over 180 cells:
5 datasets $\times$ 3 seeds $\times$ 3 models $\times$ 4 cost ratios. Datasets are our SEC arms
plus two public controls, \texttt{bank-marketing} and \texttt{creditcard}, both from OpenML.
Models are a random forest, histogram gradient boosting, and $L_2$ logistic regression.

\begin{table}[t]
\centering\small
\begin{tabular}{lrrrrrr}
\toprule
strategy & mean & median & worst & mean rank & regret & $>$floor \\
\midrule
S4 held-out validation      & 0.695 & 0.670 & 1.242 & 3.09 & 0.033 & 13\% \\
S6 analytic + OOF isotonic  & 0.765 & 0.689 & 1.968 & 4.61 & 0.103 & 23\% \\
S1 out-of-fold (5-fold)     & 0.766 & 0.702 & 1.966 & 3.49 & 0.105 & 23\% \\
S3 OOF, bootstrap median    & 0.768 & 0.682 & 1.966 & 4.08 & 0.106 & 23\% \\
S2 OOF, 3 repeats averaged  & 0.778 & 0.697 & 2.200 & 4.03 & 0.117 & 23\% \\
S7 firing-rate match        & 0.809 & 0.696 & 2.290 & 4.83 & 0.147 & 25\% \\
S5 analytic $1/(1+c)$, raw  & 0.851 & 0.696 & 3.551 & 5.14 & 0.189 & 23\% \\
S0 in-sample                & 1.156 & 0.910 & 4.581 & 6.72 & 0.494 & 39\% \\
\bottomrule
\end{tabular}
\caption{Eight threshold strategies, 180 cells, cost over the best floor. \emph{regret} is the
mean gap to an oracle threshold chosen on test, which is the only column that isolates the line
from the model. \emph{$>$floor} is how often the strategy did worse than flagging everything.}
\label{tab:strategies}
\end{table}

In-sample selection is dominated on every column. Its worst case is $4.581$, its regret is
$4.7\times$ the next worst strategy, and it is worse than doing nothing in $39\%$ of cells.

\subsection{It is worst for the models you would deploy}

The natural objection is that a weak model produced this. The opposite holds.

\begin{table}[t]
\centering\small
\begin{tabular}{llrrrr}
\toprule
dataset & model & $p/n$ & train AUC & inflation @ $c{=}20$ & val/floor \\
\midrule
bank-marketing & $L_2$ logistic  & 0.002 & 0.909 & 1.00--1.02 & 0.36 \\
bank-marketing & random forest   & 0.002 & 1.000 & \textbf{4.90--5.51} & 0.29 \\
SEC narrative  & gradient boost  & 1.385 & 1.000 & 2.62--3.76 & 0.90 \\
SEC embedding  & $L_2$ logistic  & 0.354 & 0.908 & 0.96--1.14 & 0.96 \\
\bottomrule
\end{tabular}
\caption{Inflation is the cost of a train-selected threshold over a validation-selected one, on
identical test rows, three seeds. Rows one and two are the same dataset, the same 22{,}605
training rows, the same 51 features and the same splits.}
\label{tab:strong}
\end{table}

Across all cells, every model reaching a training AUC of exactly $1.000$ shows the artefact
(inflation $1.27$ to $5.51$); no model at $0.91$ does ($1.00$ to $1.14$). The governing variable
is \textbf{train-set memorisation}, not $p/n$ and not model quality. Our own first analysis
attributed it to $p/n$; that was an artefact of having only one learner, and Table~\ref{tab:strong}
refutes it.

\subsection{Mechanism: firing-rate drift}

A threshold is chosen so that some fraction of training rows fires. Applied to new rows it fires
on a different fraction, and that drift is what the cost pays for.

\begin{center}\small
\begin{tabular}{lrrr}
\toprule
model (bank, $c{=}20$) & fires on train & fires on test & inflation \\
\midrule
random forest    & 0.117 & 0.050 & 4.90 \\
$L_2$ logistic   & 0.410 & 0.417 & 1.00 \\
\bottomrule
\end{tabular}
\end{center}

The forest's threshold was set to fire on one case in nine and fires on one in twenty. AUC gap
does \emph{not} predict the magnitude and drift does: AUC gap is a mismatch in \emph{ranking},
while a threshold is exposed to a mismatch in the score \emph{distribution} at the operating
point. This is why an uncalibrated forest, whose training scores pile at 0 and 1, is the worst
offender in the grid despite ranking well.

\subsection{Calibration cannot fix it, and this is a proof rather than a measurement}

The obvious remedy is calibration. It does nothing, in all 84 cells we ran, identically to four
decimal places.

A calibrator is a monotone map. For strictly increasing $g$, $\{s \ge t\} = \{g(s) \ge g(t)\}$,
so the family of achievable confusion matrices is unchanged and so is its cost-optimal member.
Isotonic regression is non-decreasing rather than strictly increasing, so ties could in principle
break differently; across 84 cells they never did.

Calibration fixes what a score \emph{means}. This is a problem about which rows a score
\emph{selects}.

\paragraph{Where calibration does earn its keep.} The analytic Bayes rule for this loss is to
fire when $p > 1/(1+c)$, since firing costs $(1-p)$ and not firing costs $p\,c$. That rule reads
the score's \emph{value} rather than its rank, so a calibrator is exactly what it needs. Applied
to out-of-fold isotonic scores it reproduces the empirical out-of-fold sweep in all 180 cells
(S6 versus S1 in Table~\ref{tab:strategies}: $0.765$ against $0.766$, worst case $1.968$ against
$1.966$). One tool, inert where the procedure reads ranks and load-bearing where it reads values.

Applied to \emph{raw} model output it fails on saturation, not on miscalibration: boosting's
firing rate moves only from $0.289$ to $0.340$ as $c$ goes from 3 to 20, when the sweep wants
$0.637$. Scores piled at 0 and 1 mean that dragging the threshold to $0.048$ crosses almost no
rows, so the rule cannot respond to $c$ at all.

\subsection{What does work is a size rule}

\begin{table}[t]
\centering\small
\begin{tabular}{lrrrr}
\toprule
dataset & train $n$ & in-sample & out-of-fold & validation \\
\midrule
bank-marketing        & 22{,}605 & 0.672 & \textbf{0.399} & 0.399 \\
bank-marketing (cap)  & 300      & 1.201 & 0.699 & 0.656 \\
SEC XBRL              & 2{,}167  & 1.178 & 0.874 & \textbf{0.772} \\
SEC narrative         & 2{,}167  & 1.398 & 0.912 & \textbf{0.795} \\
\bottomrule
\end{tabular}
\caption{Mean cost over floor. At 22{,}605 training rows out-of-fold is free and exactly as good
as spending a quarter of the data. At 2{,}167 it is not, and the tail is louder than the mean:
worst case $1.966$ against $1.123$.}
\label{tab:size}
\end{table}

The residual gap at small $n$ is a \emph{bias}: out-of-fold models are fitted on $(K{-}1)/K$ of an
already small set, so their scores separate less and the threshold lands low. Two natural
refinements that attack \emph{variance} therefore buy nothing, which we verified: averaging
out-of-fold scores over three repeats and taking a bootstrap-median threshold both move the mean
by less than $0.03$ and the repeated variant is marginally worse in the tail. We record this so
that others do not spend the afternoon.

\section{A checklist}

\begin{enumerate}[label=\textbf{\arabic*.}]
\item \textbf{Name the seat before naming the ratio.} If you cannot say who acts on the
prediction and what they do differently, you do not have a cost ratio; you have a missing problem
statement.
\item \textbf{Derive $c$ from hours, and report the grid.} Both halves are answerable. A point
estimate hides the uncertainty that decides the result.
\item \textbf{Never report a single-ratio asymmetric cost as a finding.} ``The model does not beat
the baseline'' is not a result. ``The model does not beat the baseline above $5.75{:}1$'' is a
different one, and a reader can act on it.
\item \textbf{Check whether your baseline is purchasable.} If always-fire is in the choice set at
$100\%$ volume with no capacity charge, your benchmark is measuring the base rate.
\item \textbf{Never select a threshold on in-sample scores.} Up to $5.5\times$, worst for your
best model, invisible in every metric you are already reporting.
\item \textbf{Out-of-fold when $n$ is large, held-out when $n$ is small.} Skip the variance
refinements; the residual is bias.
\item \textbf{Do not reach for calibration to fix a threshold.} It is provably inert unless your
rule reads values rather than ranks.
\item \textbf{Where the seat is ambiguous, report $\mathrm{precision}@k$ at a stated capacity.} It
needs no cost ratio and it is the decision the operator actually faces.
\end{enumerate}

\section{Limitations}

Three seeds on most cells and eight on two; a temporal split is reported separately and the
company split is the primary one here. The hours derivation rests on the auditor's side of the
cost, which is published, and treats the company's internal effort as unmeasured rather than
zero, so Table~\ref{tab:hours} is a lower bound on $H_\text{miss}$ and its detection-rate axis is
a stated guess rather than a measurement. Two of the three model families are tree ensembles, so
the saturation finding may not transfer to deep models whose score distributions differ. The
oracle threshold fits one free parameter to the test set and is reported as a bound rather than
an achievable number.

\section{Related work}

Cost-sensitive learning theory \citep{elkan2001} and ROC iso-performance analysis
\citep{provost2001} give the geometry we re-derive in Section~\ref{sec:sweep}; our contribution
there is the empirical crossover and the observation that it sits inside the range of asserted
constants. Instance-dependent cost \citep{bahnsen2014} motivated our firm-size weighting
experiment, which was a null. \citet{consequentialist2025} argue for $\mathrm{precision}@k$ under
known capacity, which we adopt and to which we add a variance argument: a fixed budget has no
threshold to estimate. Threshold selection is treated as a tuning detail in most applied work; we
are not aware of prior work quantifying the in-sample selection penalty as a function of
memorisation.

\bibliographystyle{plainnat}
\begin{thebibliography}{9}
\small
\bibitem[Elkan(2001)]{elkan2001}
C.~Elkan. The foundations of cost-sensitive learning. \emph{IJCAI}, 2001.
\bibitem[Provost and Fawcett(2001)]{provost2001}
F.~Provost and T.~Fawcett. Robust classification for imprecise environments.
\emph{Machine Learning}, 42(3):203--231, 2001.
\bibitem[Bahnsen et al.(2014)]{bahnsen2014}
A.~C. Bahnsen, D.~Aouada, and B.~Ottersten. Example-dependent cost-sensitive logistic regression
for credit scoring. \emph{ICMLA}, 2014.
\bibitem[GAO(2002)]{gao2002}
U.S. General Accounting Office. Financial statement restatements: trends, market impacts,
regulatory responses, and remaining challenges. GAO-03-138, 2002.
\bibitem[Interim restatements(2025)]{interim2025}
Interim restatements and the audit engagement. \emph{Journal of Accounting and Public Policy},
2025.
\bibitem[Consequentialist critique(2025)]{consequentialist2025}
A consequentialist critique of binary classification evaluation: theory, practice, and tools.
arXiv:2504.04528, 2025.
\end{thebibliography}

\end{document}
